\documentclass[10pt, aspectratio=169]{beamer}

\usepackage[T1]{fontenc}

\usepackage[UTF8]{ctex}
\usetheme[
  titlepage=moloch,
  sectionpage=progressbar,
  progressbar=frametitle,
]{moloch}

\setbeamertemplate{page number in head/foot}[appendixframenumber]
\setbeamertemplate{section in toc}[sections numbered]

\usepackage{booktabs}
\usepackage[scale=2]{ccicons}

\usepackage[semibold,light]{FiraSans}
\usepackage{FiraMono}

\usepackage{xspace}

\title{以π0为例}
\subtitle{近似计算在VLA中的应用与优化}
\date{\today}
\author{傅逸凡}
\institute{电气工程学院, 西南交通大学}
\titlegraphic{\IfFileExists{assets/logo.pdf}{\hfill
\includegraphics[width=3.5cm]{assets/logo.pdf}}{}}

\begin{document}

\maketitle

\begin{frame}
  \frametitle{Table of Contents}
  \tableofcontents
\end{frame}

\section{Introduction}

\begin{frame}
  \frametitle{研究切入点}

  \begin{columns}[T,onlytextwidth]
    \column{0.5 \textwidth}

    \begin{block}{$\pi_0$ 的模型特点}
      \small
      在视觉--语言模型的基础上，$\pi_0$ 进一步面向机器人控制任务，
      其核心区别在于：
      \begin{itemize}
        \item Prefix-Suffix 推理模式；
              % 图像语言 prefix 主要由 VLM backbone 处理并可缓存；状态和动作 suffix 由 action expert 在 flow matching 中反复计算，因此 FPGA 优化重点应放在 suffix/action expert 中重复出现的 Linear、Attention 和 MLP。
        \item Flow Matching 多步迭代推理；
              % π0 不是一次 forward 直接输出结果，而是通过 10 步左右的 denoising / Euler 更新生成动作。
              % 近似计算收益会被重复调用放大，flow steps 数量本身也可成为优化变量。
        \item 连续动作 token 带来的数值精度敏感性；
              % 输出是连续动作向量，不是离散文本 token，因此量化误差不能只看分类准确率或 top-k，而要看向量误差、cosine similarity、轨迹偏差等。
              % → 误差评估指标需要从 NLP 指标转向数值/控制相关指标。
      \end{itemize}
    \end{block}

    \begin{alertblock}{研究重点：从架构瓶颈到近似计算}
      \small
      本研究包含如下内容：

      \begin{itemize}
        \item \textbf{$\pi_0$建模、理论推导：} 分析部署瓶颈、给出近似计算优化方案；
        \item \textbf{$PyTorch$实验：} 通过量化和神经网络替代方法验证优化方案的有效性；
        \item \textbf{$Vitis HLS$测试：} 实现关键模块，评估资源占用、性能提升和误差影响；
      \end{itemize}

    \end{alertblock}


    \column{0.43\textwidth}

    \begin{exampleblock}{原题（概括）}
      \footnotesize

      \vspace{6pt}
      \textbf{任务一：分析 $\pi_0$ 推理架构与 FPGA 部署瓶颈。} \\
      说明推理阶段的数据流：图像经 ViT 等视觉编码器提取特征，再由投影网络映射到语言模型嵌入空间，作为视觉 token 参与后续生成。\\
      \vspace{2pt}
      结合 FPGA 的 LUT、BRAM、DSP 等资源，分析矩阵乘法与注意力计算复杂度高、模型参数和中间激活占用存储大、片外访存带宽压力大，以及定点化精度控制困难等问题。

      \vspace{6pt}
      \textbf{任务二：设计一种面向 $\pi_0$ 的近似计算优化方案。} \\
      量化方向：对视觉编码器、投影层和语言模型中的 Linear、Attention、FFN 等模块进行 INT8/INT4 或混合精度量化，利用 FPGA DSP 定点乘加加速；\\
      \vspace{2pt}
      神经网络替代方向：用分段线性函数、查表法或低秩近似替代复杂非线性激活、高维矩阵运算或部分注意力计算，并评估误差与硬件资源收益。
    \end{exampleblock}

  \end{columns}
\end{frame}

\begin{frame}
  \frametitle{$\pi_0$模型数据流}
  \centering
  \IfFileExists{assets/pi0_vla_flow_model.pdf}{%
    \includegraphics[width=.96\textwidth,height=.90\textheight,keepaspectratio]{assets/pi0_vla_flow_model.pdf}%
  }{%
    \IfFileExists{assets/pi0_vla_flow_model.png}{%
      \includegraphics[width=.96\textwidth,height=.90\textheight,keepaspectratio]{assets/pi0_vla_flow_model.png}%
    }{%
      \IfFileExists{pi0_vla_flow_model.png}{%
        \includegraphics[width=.96\textwidth,height=.90\textheight,keepaspectratio]{pi0_vla_flow_model.png}%
      }{%
        \begin{alertblock}{缺少架构图文件}
          请将 \texttt{pi0\_vla\_flow\_model.pdf} 或 \texttt{pi0\_vla\_flow\_model.png} 放入 \texttt{assets/} 目录。
        \end{alertblock}%
      }%
    }%
  }
  % 图中 Linear/GEMM 以 visual projector 为例，实际还包括 state/action projection、QKV/O projection、FFN 以及 action expert MLP 中的大量矩阵乘。
\end{frame}


\section{Architecture Bottleneck}

\begin{frame}
  \frametitle{任务一：$\pi_0$ 推理链路中的 FPGA 部署瓶颈}
  \small
  \begin{columns}[T,onlytextwidth]
    \column{0.48\textwidth}
    \begin{block}{为什么不做完整端到端部署}
      \begin{itemize}
        \item $\pi_0$ 不是单纯 VLM，而是 VLM backbone + action expert + flow matching。
        \item 完整模型参数量大，推理链路包含图像、语言、状态、动作和多步积分。
        \item 一周内更合理的路线是\textbf{模块级分析与验证}，而非完整 checkpoint 部署到 FPGA。
      \end{itemize}
    \end{block}

    \begin{exampleblock}{本项目降级策略}
      选择 red dashed boxes 对应的关键模块：Linear/GEMM、QKV、FFN、Softmax、GELU/SiLU、RMSNorm、Action Expert MLP 与 Flow Step Reduction。
    \end{exampleblock}

    \column{0.48\textwidth}
    \begin{alertblock}{主要瓶颈}
      \begin{tabular}{@{}p{0.25\linewidth}p{0.68\linewidth}@{}}
        \toprule
        类型 & 具体问题 \\
        \midrule
        计算量 & QKV、FFN、projector 中 GEMM 占主导 \\
        存储 & VLM 权重和中间激活远超片上 BRAM \\
        访存 & 大矩阵权重反复搬运，带宽压力明显 \\
        非线性 & Softmax、GELU、RMSNorm 含 exp/div/sqrt \\
        多步推理 & flow matching 需要多次 action expert 调用 \\
        \bottomrule
      \end{tabular}
    \end{alertblock}
  \end{columns}
\end{frame}

\section{Approximation Strategy}

\begin{frame}
  \frametitle{任务二：面向 $\pi_0$ 的模块级近似计算方案}
  \small
  \begin{table}
    \begin{tabular}{@{}p{0.24\linewidth}p{0.38\linewidth}p{0.28\linewidth}@{}}
      \toprule
      优化对象 & 近似/定点化方法 & 预期收益 \\
      \midrule
      Linear / QKV / FFN & INT8、INT4、W4A8 fake quant；INT32/INT24 累加后 requantize & 降低 GEMM 计算和权重存储 \\
      Softmax & LUT、PWL、Taylor、base-2 softmax & 避免高成本 exp/div \\
      GELU / SiLU & PWL、hard approximation、clipped linear & 降低复杂激活函数开销 \\
      RMSNorm & fixed-point + reciprocal sqrt 近似 & 降低 sqrt/div 成本 \\
      Flow steps & 10 steps $\rightarrow$ 8/6/4/2 steps & 减少 action expert 调用次数 \\
      \bottomrule
    \end{tabular}
  \end{table}

  \vspace{4pt}
  \begin{alertblock}{核心思路}
    不直接宣称完整 $\pi_0$ 可部署到 FPGA，而是把模型拆成可验证的算子与模块，用误差、延迟和存储压缩比评估近似计算可行性。
  \end{alertblock}
\end{frame}

\section{PyTorch Experiments}

\begin{frame}
  \frametitle{PyTorch 实验设计：从随机 Proxy 到真实权重 Proxy}
  \small
  \begin{columns}[T,onlytextwidth]
    \column{0.52\textwidth}
    \begin{block}{实验分层}
      \begin{tabular}{@{}p{0.30\linewidth}p{0.28\linewidth}p{0.30\linewidth}@{}}
        \toprule
        实验层级 & 权重/输入 & 目的 \\
        \midrule
        random + $\pi_0$ shape & 随机权重 + 随机输入 & 验证同结构同规模趋势 \\
        real weight proxy & $\pi_0$ 权重 + 随机输入 & 验证真实参数分布下误差 \\
        toy flow & 合成动作轨迹 & 验证步数减少的误差--延迟权衡 \\
        \bottomrule
      \end{tabular}
    \end{block}

    \begin{exampleblock}{评价指标}
      MSE、MAE、max error、cosine similarity、relative L2、KL divergence、latency、weight size、compression ratio。
    \end{exampleblock}

    \column{0.44\textwidth}
    \begin{alertblock}{主动限定结论边界}
      \begin{itemize}
        \item 不是完整 $\pi_0$ 端到端量化评估；
        \item 没有真实 observation / activation / action dataset；
        \item 不能推出真实机器人任务成功率下降；
        \item 结论只支撑\textbf{模块级 proxy benchmark}。
      \end{itemize}
    \end{alertblock}
  \end{columns}
\end{frame}

\begin{frame}
  \frametitle{实验覆盖的 $\pi_0$ 关键模块与 Shape}
  \scriptsize
  \begin{table}
    \begin{tabular}{@{}p{0.25\linewidth}p{0.37\linewidth}p{0.30\linewidth}@{}}
      \toprule
      模块 & 实验 shape / 结构 & 对应红框 \\
      \midrule
      Visual projector & $[1,768,1152] \rightarrow [1,768,2048]$ & Linear / GEMM \\
      VLM attention & $q:2048\rightarrow2048$, $k/v:2048\rightarrow256$, $o:2048\rightarrow2048$ & QKV Linear, Softmax \\
      VLM FFN & $2048 \rightarrow 16384 \rightarrow 2048$，gated FFN & FFN, GELU/SiLU \\
      Action expert FFN & $1024 \rightarrow 4096 \rightarrow 1024$ & Action Expert MLP \\
      State/action projection & $32\rightarrow1024$, $1024\rightarrow32$ & Linear / GEMM \\
      RMSNorm & $[1,768,2048]$, $[1,50,1024]$ & RMSNorm \\
      Flow matching toy & $H=50$, action dim $=32$, steps $=10/8/6/4/2$ & Flow Step Reduction \\
      \bottomrule
    \end{tabular}
  \end{table}

  \begin{block}{实验定位}
    先用随机张量对齐结构和规模，再用 $\pi_0$ checkpoint 中提取的真实模块权重重复测试，提升 proxy benchmark 的可信度。
  \end{block}
\end{frame}

\begin{frame}
  \frametitle{实验结果一：GEMM 类模块适合优先 INT8 定点化}
  \small
  \begin{columns}[T,onlytextwidth]
    \column{0.56\textwidth}
    \begin{table}
      \begin{tabular}{@{}p{0.40\linewidth}cc@{}}
        \toprule
        实验设置 & INT8 最低 cosine & INT4 最低 cosine \\
        \midrule
        $\pi_0$-aligned random & 0.998402 & 0.963150 \\
        real $\pi_0$ weight & 0.994843 & 0.921080 \\
        \bottomrule
      \end{tabular}
    \end{table}

    \begin{exampleblock}{观察}
      \begin{itemize}
        \item INT8 在随机 proxy 和真实权重 proxy 中都较稳定；
        \item INT4/W4A8 权重压缩更强，但误差风险明显增加；
        \item GEMM 类模块是 FPGA 定点化最优先对象。
      \end{itemize}
    \end{exampleblock}

    \column{0.40\textwidth}
    \begin{alertblock}{对应 $\pi_0$ 模块}
      \begin{itemize}
        \item Visual projector；
        \item QKV/O projection；
        \item FFN gate/up/down；
        \item State/action projection；
        \item Action expert MLP。
      \end{itemize}
    \end{alertblock}
  \end{columns}
\end{frame}

\begin{frame}
  \frametitle{实验结果二：Softmax / RMSNorm 可近似，激活替换需谨慎}
  \small
  \begin{columns}[T,onlytextwidth]
    \column{0.52\textwidth}
    \begin{block}{Softmax}
      真实 Q/K 权重生成 attention score 后，LUT softmax 的 relative L2 约为 $4.3\times10^{-5}$，说明 LUT 是较好的硬件友好候选。
    \end{block}

    \begin{block}{RMSNorm}
      真实权重 proxy 中 RMSNorm 近似最大 relative L2 约为 0.0117，reciprocal sqrt 近似相对可控。
    \end{block}

    \column{0.44\textwidth}
    \begin{alertblock}{GELU / SiLU}
      FFN 激活替换必须放回完整 FFN 评价。identity 或过激 clipped-linear 替换可能导致 relative L2 接近 1，不能直接作为默认方案。
    \end{alertblock}

    \begin{exampleblock}{结论}
      非线性算子中，Softmax 与 RMSNorm 的近似更稳；GELU/SiLU 应优先采用保守 PWL/LUT，而非粗暴线性替换。
    \end{exampleblock}
  \end{columns}
\end{frame}

\begin{frame}
  \frametitle{实验结果三：Flow Step Reduction 是算法级加速旋钮}
  \small
  \begin{columns}[T,onlytextwidth]
    \column{0.50\textwidth}
    \begin{block}{实验设置}
      \begin{itemize}
        \item toy flow matching，不运行完整 $\pi_0$；
        \item 对齐 action shape：$H=50$, action dim $=32$；
        \item 对比 10 / 8 / 6 / 4 / 2 step Euler integration。
      \end{itemize}
    \end{block}

    \column{0.46\textwidth}
    \begin{exampleblock}{观察}
      2-step 相比 10-step 约 $4.46\times$ 加速，MSE 约 $4.32\times10^{-3}$，cosine 约 0.9975。
    \end{exampleblock}

    \begin{alertblock}{边界}
      该结果只能说明步数减少存在误差--延迟权衡，不能代表真实机器人动作质量。
    \end{alertblock}
  \end{columns}
\end{frame}

\section{Conclusion}

\begin{frame}
  \frametitle{总结：$\pi_0$ 模块级近似计算验证结论}
  \small
  \begin{enumerate}
    \item \textbf{架构层面：} $\pi_0$ 是 VLM backbone + action expert + flow matching，不适合直接完整 FPGA 部署。
    \item \textbf{GEMM 层面：} visual projector、QKV/O、FFN、state/action projection 和 action expert MLP 均适合优先 INT8 定点化。
    \item \textbf{非线性层面：} Softmax 可采用 LUT/PWL，RMSNorm 可采用 reciprocal sqrt 近似；GELU/SiLU 需要保守处理。
    \item \textbf{算法层面：} Flow step reduction 有延迟收益，但必须后续结合真实 action trajectory 验证。
    \item \textbf{工程层面：} 本工作已完成 PyTorch 模块级验证，后续可将 INT8 tiled GEMM、LUT softmax、PWL GELU 落到 Vitis HLS kernel。
  \end{enumerate}
\end{frame}

\begin{frame}[standout]
  谢谢老师，请批评指正。
\end{frame}

\appendix

\section{Appendix}

\begin{frame}
  \frametitle{Appendix：真实权重 Proxy 的作用与限制}
  \small
  \begin{columns}[T,onlytextwidth]
    \column{0.48\textwidth}
    \begin{block}{作用}
      \begin{itemize}
        \item 使用 $\pi_0$ checkpoint 中真实模块权重；
        \item 反映真实参数分布下的量化误差；
        \item 比纯随机权重更接近实际模型。
      \end{itemize}
    \end{block}

    \column{0.48\textwidth}
    \begin{alertblock}{限制}
      \begin{itemize}
        \item 输入仍是 random tensor；
        \item 没有真实 activation calibration；
        \item 没有 action trajectory / success rate；
        \item 因此不等于端到端 $\pi_0$ 量化评估。
      \end{itemize}
    \end{alertblock}
  \end{columns}
\end{frame}

\begin{frame}
  \frametitle{Appendix：后续 Vitis HLS 可做的最小实验}
  \small
  \begin{table}
    \begin{tabular}{@{}p{0.24\linewidth}p{0.38\linewidth}p{0.28\linewidth}@{}}
      \toprule
      HLS kernel & 设计要点 & 对应模块 \\
      \midrule
      INT8 tiled GEMM & loop tiling、pipeline、unroll、array partition & projector、QKV、FFN、action projection \\
      LUT softmax & fixed-point input、LUT exp、reciprocal 近似 & attention softmax \\
      PWL GELU & 分段线性拟合、查表参数 & FFN / action expert MLP \\
      RMSNorm & fixed-point square/mean + rsqrt 近似 & Transformer / action expert \\
      \bottomrule
    \end{tabular}
  \end{table}
\end{frame}

\end{document}
